\documentclass[UTF8,a4paper,11pt]{ctexart}

\usepackage[margin=2.4cm]{geometry}
\usepackage{amsmath}
\usepackage{booktabs}
\usepackage{longtable}
\usepackage{array}
\usepackage{xcolor}
\usepackage{hyperref}
\usepackage{enumitem}

\hypersetup{
  colorlinks=true,
  linkcolor=blue,
  urlcolor=blue,
  citecolor=blue
}

\setlist[itemize]{leftmargin=2em}
\setlist[enumerate]{leftmargin=2em}
\renewcommand{\arraystretch}{1.15}
\emergencystretch=3em

\title{\textbf{LLMRepair Pipeline 说明}\\
\large 基于 LLM 语义 Profiling 和 RNNRepair 式状态抽象的 Who\&When 失败归因}
\author{Model-based LLM Agent Debugging}
\date{\today}

\begin{document}
\maketitle

\section{我们要解决什么问题}

Who\&When 数据集里的每一条样本是一段失败的多智能体执行日志。日志里有一个任务问题、多个 agent 的历史对话或工具调用，以及人工标注的两个答案：

\begin{itemize}
  \item \textbf{Who}: 哪个 agent 对最终失败负主要责任，即 \texttt{mistake\_agent}。
  \item \textbf{When}: 决定性错误最早发生在哪一步，即 \texttt{mistake\_step}。
\end{itemize}

我们的目标不是让 LLM 直接 all-at-once 输出 ``谁错了、哪一步错了''。我们的目标是做一个更结构化的方法：

\begin{center}
\fbox{\parbox{0.88\linewidth}{
\centering
LLM 先把每一步变成结构化语义标签；\\
再用类似 RNNRepair 的状态抽象方法，把这些标签变成 agent 状态机；\\
最后根据训练集中哪些状态更容易出错，定位新日志里的 Who 和 When。
}}
\end{center}

所以这个方法仍然是 LLM-based，因为每一步的语义理解由 DeepSeek-v4-flash 完成；但最终定位不是一次性问 LLM，而是由一个可解释、可缓存、可交叉验证的状态模型完成。

\section{当前已经跑通的实验配置}

当前已经完成的是 \textbf{Hand-Crafted} 子集上的 profiling 和五折评估。

\begin{table}[htbp]
\centering
\begin{tabular}{ll}
\toprule
项目 & 当前配置 \\
\midrule
数据集 & Who\&When Hand-Crafted \\
日志数量 & 58 条 \\
总 step 数 & 2993 \\
非 human step 数 & 2935 \\
LLM & DeepSeek-v4-flash \\
是否给 ground truth answer & 否，\texttt{without\_ground\_truth} \\
每步局部上下文窗口 & 当前 step 之前最多 4 步 \\
评估方式 & log-level 5-fold cross validation \\
状态模型 & 每个 agent 单独训练 PCA + GMM \\
\bottomrule
\end{tabular}
\end{table}

非 human step 的 agent 分布如下：

\begin{table}[htbp]
\centering
\begin{tabular}{lr}
\toprule
Agent & step 数 \\
\midrule
Orchestrator & 2283 \\
WebSurfer & 568 \\
Assistant & 39 \\
FileSurfer & 35 \\
ComputerTerminal & 10 \\
\bottomrule
\end{tabular}
\end{table}

\section{整体流程}

整个 pipeline 可以分成六步：

\begin{enumerate}
  \item 把 Who\&When JSON 日志摊平成 ``每一步一行'' 的表。
  \item 对每一个非 human step 调用 DeepSeek-v4-flash，得到结构化语义 profile。
  \item 把语义 profile 转成数值特征。
  \item 每个 agent 单独训练 GMM，把它的 step 聚成若干抽象状态。
  \item 在训练集上统计每种状态或转移的历史错误风险，得到 risk table。
  \item 在测试日志上逐步打分，预测 responsible agent 和 decisive error step。
\end{enumerate}

一句话概括就是：

\begin{center}
\fbox{\parbox{0.88\linewidth}{
\centering
原始日志 $\rightarrow$ 每步语义 profile $\rightarrow$ agent 状态机 $\rightarrow$ 状态风险表 $\rightarrow$ Who/When 预测
}}
\end{center}

\section{第一步：把日志摊平成每步一行}

原始 JSON 里，一条日志大致有这些内容：

\begin{itemize}
  \item \texttt{question}: 任务问题。
  \item \texttt{history}: agent 交互历史，是一个 step 列表。
  \item \texttt{mistake\_agent}: 人工标注的 responsible agent。
  \item \texttt{mistake\_step}: 人工标注的 decisive error step。
\end{itemize}

我们先把每条日志的 \texttt{history} 摊平。摊平之后，每一行代表一个 step，核心字段如下：

\begin{longtable}{p{0.28\linewidth}p{0.62\linewidth}}
\toprule
字段 & 含义 \\
\midrule
\endhead
\texttt{log\_id} & 属于哪一条日志。 \\
\texttt{step\_idx} & 这是日志里的第几步，从 0 开始。 \\
\texttt{agent\_name} & 归一化后的 agent 名字，例如 \texttt{Orchestrator}, \texttt{WebSurfer}。 \\
\texttt{raw\_agent} & 原始 JSON 里的 agent 名字。 \\
\texttt{content} & 这一条 step 的实际文本或工具结果。 \\
\texttt{question} & 当前日志对应的任务问题。 \\
\texttt{ground\_truth} & 如果数据里有最终正确答案，会保留；当前实验没有喂给 LLM。 \\
\texttt{mistake\_agent} & 人工标注的错误 agent。 \\
\texttt{mistake\_step} & 人工标注的错误 step。 \\
\texttt{is\_decisive\_error} & 如果 \texttt{step\_idx == mistake\_step}，则为 true。 \\
\bottomrule
\end{longtable}

这一步输出到：

\begin{verbatim}
artifacts/handcrafted_steps.jsonl
\end{verbatim}

对应命令是：

\begin{verbatim}
python scripts/llmrepair_pipeline.py prepare-steps \
  --dataset-root datasets/Who_and_When \
  --output artifacts/handcrafted_steps.jsonl
\end{verbatim}

\section{第二步：LLM semantic profiling}

这是整个方法里最关键的 LLM-based 部分。对每一个非 human step，我们调用一次 DeepSeek-v4-flash。
但是我们\textbf{不}让它直接回答 ``谁错了、哪一步错了''。我们只让它判断当前 step 的局部语义状态。

\subsection{每次 LLM 调用看什么}

对于某个 step $t$，LLM 看到的是：

\begin{enumerate}
  \item 当前日志的任务问题 \texttt{question}。
  \item 当前 step 前面最多 4 步局部上下文。
  \item 当前 step 的 \texttt{step\_idx}。
  \item 当前 step 的 agent 名字。
  \item 当前 step 的内容 \texttt{content}。
\end{enumerate}

当前实验是 \texttt{without\_ground\_truth}，所以 prompt 里没有给最终正确答案。

注意，这里不是把完整 history 都塞给 LLM。当前实现采用局部窗口：

\[
\text{context}(t) = \{t-4, t-3, t-2, t-1\}
\]

如果前面不足 4 步，就给已有的前文。选择 4 不是理论常数，而是当前工程默认值：它能提供足够的局部因果上下文，同时控制 token 成本和噪声。后续可以把窗口 0、2、4、8 做 ablation。

\subsection{LLM 输出什么}

LLM 必须返回一个 JSON 对象。当前字段是：

\begin{longtable}{p{0.30\linewidth}p{0.58\linewidth}}
\toprule
字段 & 含义 \\
\midrule
\endhead
\texttt{action\_type} & 当前 step 的动作类型，如 plan, delegate, web\_search, web\_browse, tool\_result, file\_read, code\_execute, reason, answer, termination, other。 \\
\texttt{agent\_intent} & agent 这一步想做什么，短文本。 \\
\texttt{task\_alignment} & 0 到 3 的整数，越高表示越贴近任务。 \\
\texttt{evidence\_quality} & 0 到 3 的整数，越高表示证据越强。 \\
\texttt{progress\_state} & on\_track, uncertain, stalled, off\_track 之一。 \\
\texttt{new\_claim} & 是否提出了新结论或新声明。 \\
\texttt{unsupported\_claim} & 是否有缺少证据支持的声明。 \\
\texttt{tool\_use} & 是否涉及工具使用。 \\
\texttt{tool\_error} & 工具是否失败或返回异常。 \\
\texttt{reasoning\_error\_signal} & 是否出现明显推理错误信号。 \\
\texttt{state\_summary} & 对当前局部状态的一句话总结。 \\
\bottomrule
\end{longtable}

可以把这一步理解成：LLM 是一个 ``语义传感器''。它不做最终判决，只把自然语言日志翻译成结构化观察。

\subsection{为什么这样设计}

直接 all-at-once 问 LLM 有几个问题：

\begin{itemize}
  \item 长日志容易超上下文或注意力分散。
  \item LLM 的最终答案不稳定，输出格式也可能不稳定。
  \item 很难知道它到底为什么选某一步。
\end{itemize}

我们现在把任务拆开：

\begin{itemize}
  \item LLM 负责它擅长的事：理解局部文本语义。
  \item 状态模型负责可重复的事：聚类、统计风险、交叉验证评估。
\end{itemize}

这使得结果更可缓存、更可解释，也更接近 RNNRepair 的思想。

\section{第三步：把 profile 转成数值特征}

LLM 输出的 JSON 不能直接喂给 GMM，所以要把它变成特征向量。

当前真正进入状态模型的字段是：

\begin{itemize}
  \item 数值或布尔特征：\texttt{task\_alignment}, \texttt{evidence\_quality}, \texttt{new\_claim}, \texttt{unsupported\_claim}, \texttt{tool\_use}, \texttt{tool\_error}, \texttt{reasoning\_error\_signal}。
  \item 类别特征：\texttt{action\_type}, \texttt{progress\_state}。
\end{itemize}

类别特征用 one-hot 编码。例如 \texttt{action\_type=web\_search} 会变成一个二进制维度。

\texttt{agent\_intent} 和 \texttt{state\_summary} 当前主要用于人工解释，暂时没有进入 GMM。这样做是为了先保持 v1 简单、稳定。如果后面需要，可以把它们再 embedding 成向量加入特征。

\section{第四步：每个 agent 单独建状态机}

RNNRepair 的核心思想是：不要直接看原始高维 hidden state，而是先把 hidden state 聚成少数几个抽象状态，然后分析哪些状态和错误有关。

我们的对象不是 RNN hidden state，而是 LLM semantic profile。因此对应关系是：

\begin{table}[htbp]
\centering
\begin{tabular}{lll}
\toprule
RNNRepair & 我们的方法 & 解释 \\
\midrule
RNN hidden state & LLM semantic feature & 每一步的内部状态表示 \\
GMM abstract state & GMM semantic state & 把相似 step 聚成状态 \\
state confidence/stability & semantic stability & 判断状态语义是否纯 \\
influence model & risk table & 查某种状态历史上多危险 \\
\bottomrule
\end{tabular}
\end{table}

\subsection{为什么每个 agent 单独建}

不同 agent 的行为空间不一样。比如：

\begin{itemize}
  \item WebSurfer 的状态可能是搜索、浏览、证据不足、网页错误。
  \item Orchestrator 的状态可能是规划、分派、总结、错误终止。
  \item ComputerTerminal 的状态可能是执行代码、工具输出、命令失败。
\end{itemize}

如果把所有 agent 混在一起聚类，状态会混乱。因此当前做法是：

\[
\text{for each agent } a,\quad \text{fit a separate GMM on steps spoken by } a
\]

\subsection{训练 GMM 前的处理}

对每个 agent 的所有训练 step：

\begin{enumerate}
  \item 把 profile 字段转成数值向量。
  \item 用 \texttt{StandardScaler} 标准化。
  \item 用 PCA 降维，当前最多降到 8 维。
  \item 用 GMM 聚类得到抽象状态。
\end{enumerate}

这里 PCA 的作用是让 GMM 更稳定。原始 one-hot 特征维度不算特别高，但数据量很小，尤其 Assistant、FileSurfer、ComputerTerminal 的样本很少，直接 GMM 容易不稳。

\section{第五步：GMM 的 K 怎么选}

这里我们参考 RNNRepair 的 ``semantic-guided abstraction'' 思路，不是只看 BIC，也不是用错误标签直接调 K。

对候选 $K=2,3,\ldots,15$，我们都训练一个 GMM。然后看聚出来的每个状态语义是否稳定。

\subsection{什么叫状态语义稳定}

假设某个状态 $q$ 里有很多 step。我们看这些 step 的语义字段是否一致。例如：

\begin{itemize}
  \item 这个状态里的 \texttt{progress\_state} 大多数是不是都是 \texttt{off\_track}？
  \item 这个状态里的 \texttt{action\_type} 大多数是不是都是 \texttt{web\_browse}？
  \item 这个状态里的 \texttt{unsupported\_claim} 大多数是不是都是 true？
\end{itemize}

对每个语义字段，我们计算纯度：

\[
\text{purity}(q, c) =
\frac{\text{状态 }q\text{ 中字段 }c\text{ 的最大类别计数}}
{\text{状态 }q\text{ 的样本数}}
\]

然后对所有语义字段取平均，得到这个状态的 stability：

\[
\text{stability}(q) =
\frac{1}{|C|}\sum_{c\in C}\text{purity}(q,c)
\]

再对所有状态取平均：

\[
\text{avg\_stability}(K) =
\frac{1}{K}\sum_{q=1}^{K}\text{stability}(q)
\]

\subsection{当前选择规则}

当前规则是：

\begin{center}
\fbox{\parbox{0.86\linewidth}{
从小到大尝试 K，选择第一个满足
\[
\text{avg\_stability}(K)\geq 0.72
\]
并且每个状态至少有 5 个样本的 K。
如果没有 K 同时满足，就选择 stability 最好的可用 K。
}}
\end{center}

这和 RNNRepair 的精神一致：我们希望状态不是为了贴标签而过拟合，而是语义上真的能解释。

在全量 hand-crafted profile 上，当前选出来的 K 是：

\begin{table}[htbp]
\centering
\begin{tabular}{lrl}
\toprule
Agent & Selected K & 选择原因 \\
\midrule
Orchestrator & 2 & 达到 stability 阈值 \\
WebSurfer & 3 & 达到 stability 阈值 \\
Assistant & 2 & 达到 stability 阈值 \\
FileSurfer & 2 & 达到 stability 阈值 \\
ComputerTerminal & 2 & 小样本下的 best available \\
\bottomrule
\end{tabular}
\end{table}

\section{第六步：训练 risk table}

有了每一步的 agent 状态之后，我们要统计：训练集中哪些状态更容易成为决定性错误。

首先定义全局基础错误率：

\[
p_0 = \frac{\text{训练集中决定性错误 step 数}}{\text{训练集中可用 step 数}}
\]

然后对某个 key，例如某个 agent 的某个状态，计算：

\[
p(\text{error}\mid key)
=
\frac{\text{key 下错误 step 数}+\alpha}
{\text{key 下总 step 数}+2\alpha}
\]

这里使用 Laplace smoothing，当前：

\[
\alpha = 0.5
\]

最后定义：

\[
\text{risk\_lift}(key)
=
\frac{p(\text{error}\mid key)}{p_0}
\]

直观解释：

\begin{itemize}
  \item \texttt{risk\_lift = 1}: 这个状态和平均水平一样危险。
  \item \texttt{risk\_lift > 1}: 这个状态比平均更容易出错。
  \item \texttt{risk\_lift < 1}: 这个状态比平均更安全。
\end{itemize}

\subsection{risk table 的 key}

当前有四层 key，从具体到粗糙依次回退：

\begin{enumerate}
  \item \textbf{transition}: \((agent,\ previous\_state,\ action\_type,\ current\_state)\)
  \item \textbf{state}: \((agent,\ current\_state)\)
  \item \textbf{action}: \((agent,\ action\_type)\)
  \item \textbf{agent}: \((agent)\)
\end{enumerate}

预测时先查最具体的 transition。如果样本支持数不够，就回退到 state，再回退到 action，再回退到 agent。如果都没有足够支持，就使用基础风险，即 \texttt{risk\_lift = 1}。

\subsection{为什么训练时只用 mistake\_step 之前的 step}

训练 risk table 时，只使用：

\[
step\_idx \leq mistake\_step
\]

原因是：决定性错误之后的 step 往往是错误的后果，而不是错误的原因。如果把后续 step 也当作正常训练样本，模型可能会学到 ``错误之后的混乱状态''，而不是 ``错误发生的状态''。

\section{第七步：每一步怎么打分}

测试一条新日志时，我们对每个非 human step 都算三个东西：

\begin{enumerate}
  \item \texttt{profile\_score}: LLM 语义 profile 自身显示的局部可疑程度。
  \item \texttt{risk\_lift}: 这个状态在训练集中历史上有多危险。
  \item \texttt{combined\_score}: 把上面两者合起来的总分。
\end{enumerate}

\subsection{profile\_score}

当前公式是：

\[
\begin{aligned}
\text{profile\_score}
=&\ 2.0(3-\text{task\_alignment})
+0.75(3-\text{evidence\_quality})\\
&+\text{progress\_penalty}
+2.0\cdot I(\text{reasoning\_error\_signal})\\
&+1.5\cdot I(\text{tool\_error})
+1.5\cdot I(\text{unsupported\_claim})
\end{aligned}
\]

其中：

\[
\text{progress\_penalty}=
\begin{cases}
0, & \text{on\_track}\\
0.25, & \text{uncertain}\\
1.0, & \text{stalled}\\
2.0, & \text{off\_track}
\end{cases}
\]

这个分数只看 LLM 对当前 step 的语义判断。比如 task alignment 很低、evidence quality 很低、又有 unsupported claim，那么 profile score 就会高。

\subsection{combined\_score}

当前组合公式是：

\[
\text{combined\_score}
=
\text{profile\_score}
+2\log(\text{risk\_lift})
\]

这里的含义是：

\begin{itemize}
  \item \texttt{profile\_score}: 当前 step 的文本内容本身像不像错误。
  \item \texttt{risk\_lift}: 训练集中同类状态是否真的经常出错。
\end{itemize}

为什么用 \(\log\)？因为 risk lift 是倍数。如果直接加倍数，少数小样本高风险状态会把分数拉得过猛。取 log 后：

\[
\log(1)=0,\quad
\log(2)\approx 0.69,\quad
\log(4)\approx 1.39,\quad
\log(0.5)\approx -0.69
\]

所以：

\begin{itemize}
  \item \texttt{risk\_lift = 1} 时，不改变分数。
  \item \texttt{risk\_lift > 1} 时，加分。
  \item \texttt{risk\_lift < 1} 时，扣分。
  \item 特别大的 lift 会被 log 压住。
\end{itemize}

前面的 2 是当前经验权重，意思是让历史状态风险有明显影响，但不完全压过 LLM 的局部语义判断。后续论文实验里，这个权重应该做验证集调参和 ablation。

\section{第八步：怎么从 step 分数得到 Who 和 When}

当前实现了五种 decoding 模式：

\begin{longtable}{p{0.26\linewidth}p{0.62\linewidth}}
\toprule
模式 & 含义 \\
\midrule
\endhead
\texttt{profile\_argmax} & 选择 \texttt{profile\_score} 最高的 step。 \\
\texttt{profile\_earliest4} & 选择最早一个 \texttt{profile\_score >= 4} 的 step；如果没有，再退回 argmax。 \\
\texttt{risk\_argmax} & 选择 \texttt{risk\_lift} 最高的 step。 \\
\texttt{combined\_argmax} & 选择 \texttt{combined\_score} 最高的 step。 \\
\texttt{combined\_earliest4} & 选择最早一个 \texttt{combined\_score >= 4} 的 step；如果没有，再退回 combined argmax。 \\
\bottomrule
\end{longtable}

如果预测出来的 step 是第 $t$ 步，那么：

\begin{itemize}
  \item \textbf{When} 就预测为 $t$。
  \item \textbf{Who} 就预测为第 $t$ 步对应的 agent。
\end{itemize}

当前建议报告两个主结果：

\begin{itemize}
  \item \textbf{LLMRepair-Agent}: 使用 \texttt{combined\_argmax}，因为它的 agent-level accuracy 最高。
  \item \textbf{LLMRepair-Step}: 使用 \texttt{combined\_earliest4}，因为它更符合 ``最早决定性错误'' 的定义，step 容差指标更好。
\end{itemize}

\section{第九步：五折交叉验证怎么做}

我们用 5-fold cross validation，并且按 log 切分，而不是按 step 切分。

每一折：

\begin{enumerate}
  \item 58 条日志被分成训练日志和测试日志。
  \item 只用训练日志拟合每个 agent 的 PCA + GMM。
  \item 只用训练日志构建 risk table。
  \item 在测试日志上预测 Who 和 When。
  \item 收集 5 折所有测试预测，最后统一算指标。
\end{enumerate}

这样可以避免同一条日志的一部分 step 进入训练、另一部分 step 进入测试造成泄漏。

LLM semantic profile 是 label-free 的，也就是 prompt 不包含 \texttt{mistake\_agent} 或 \texttt{mistake\_step}，因此可以预先缓存。
真正使用人工错误标签的是 GMM 之后的 risk table 和评估阶段。

\section{当前结果}

当前 hand-crafted 58 条日志上的五折结果如下：

\begin{table}[htbp]
\centering
\small
\begin{tabular}{lrrrrr}
\toprule
方法 & Agent Acc. & Step Acc. & Tol@3 & Tol@5 & Avg Dist. \\
\midrule
\texttt{profile\_argmax} & 56.90\% & 10.34\% & 27.59\% & 39.66\% & 17.79 \\
\texttt{profile\_earliest4} & 53.45\% & 15.52\% & 34.48\% & 53.45\% & 9.16 \\
\texttt{risk\_argmax} & 56.90\% & 22.41\% & 31.03\% & 46.55\% & 14.43 \\
\texttt{combined\_argmax} & \textbf{70.69\%} & 18.97\% & 31.03\% & 41.38\% & 17.95 \\
\texttt{combined\_earliest4} & 58.62\% & 20.69\% & \textbf{39.66\%} & \textbf{58.62\%} & 9.29 \\
\bottomrule
\end{tabular}
\end{table}

这个表的解读是：

\begin{itemize}
  \item 如果主要看 \textbf{Who}，\texttt{combined\_argmax} 最强，agent accuracy 达到 70.69\%。
  \item 如果主要看 \textbf{When}，\texttt{risk\_argmax} 的 exact step accuracy 最高，为 22.41\%。
  \item 但 \texttt{combined\_earliest4} 的容差指标更好，Tol@5 达到 58.62\%，也更符合 ``最早错误'' 的定义。
  \item \texttt{profile\_score} 和 \texttt{risk\_lift} 结合后，确实比只用其中之一更适合 agent attribution。
\end{itemize}

\section{和论文 baseline 的关系}

论文里的 all-at-once、step-by-step、binary-search 是直接让 LLM 判断 Who 和 When。我们这里的方法不同：

\begin{itemize}
  \item 论文 baseline 是 \textbf{zero-shot LLM judge}。
  \item 我们的方法是 \textbf{LLM semantic profiler + trainable state/risk model}。
\end{itemize}

所以比较时必须说明：我们不是同一种 inference setting。我们的优势是可以从训练折学习状态风险；代价是需要标注训练日志。

根据本地整理的论文 hand-crafted 无 ground truth 数字，GPT-4o baseline 大致是：

\begin{table}[htbp]
\centering
\begin{tabular}{lrr}
\toprule
方法 & Agent Acc. & Step Acc. \\
\midrule
All-at-once & 53.44\% & 3.51\% \\
Step-by-step & 32.75\% & 8.77\% \\
Binary search & 36.21\% & 6.90\% \\
\bottomrule
\end{tabular}
\end{table}

从数值上看，我们当前结果在 hand-crafted 上明显高于这些 baseline，尤其是 agent-level 和 step exact。但论文写作时要避免夸大，必须写清楚这是 \textbf{supervised cross-validation setting}，不是纯 zero-shot。

\section{实际运行命令}

完整命令如下。这里不写 API key，运行前需要在环境变量中设置 \texttt{DEEPSEEK\_API\_KEY}。

\begin{verbatim}
# 1. 摊平 hand-crafted logs
python scripts/llmrepair_pipeline.py prepare-steps \
  --dataset-root datasets/Who_and_When \
  --output artifacts/handcrafted_steps.jsonl

# 2. 对每个非 human step 调 DeepSeek 做 semantic profiling
python scripts/llmrepair_pipeline.py profile-steps \
  --steps artifacts/handcrafted_steps.jsonl \
  --output artifacts/profiles/deepseek-v4-flash_without_gt.jsonl \
  --model deepseek-v4-flash \
  --ground-truth-mode without_ground_truth \
  --workers 4

# 3. 全量 profile 上看每个 agent 的 K 选择
python scripts/llmrepair_pipeline.py select-k \
  --profiles artifacts/profiles/deepseek-v4-flash_without_gt.jsonl \
  --output artifacts/state_models/k_selection_without_gt.json \
  --theta 0.72 \
  --min-support 5

# 4. 五折评估
python scripts/llmrepair_pipeline.py evaluate \
  --steps artifacts/handcrafted_steps.jsonl \
  --profiles artifacts/profiles/deepseek-v4-flash_without_gt.jsonl \
  --output-dir artifacts/eval/without_gt \
  --n-splits 5 \
  --random-state 0
\end{verbatim}

当前 profiling 已经完成，缓存文件里有 2935 个非 human step 的 profile。总 token 统计约为：

\begin{itemize}
  \item prompt tokens: 2,675,846
  \item completion tokens: 1,279,194
  \item total tokens: 3,955,040
\end{itemize}

\section{这个方法的主要创新点}

当前 pipeline 的创新点可以这样概括：

\begin{enumerate}
  \item \textbf{把 LLM 从最终裁判改成语义传感器。} LLM 不直接输出 Who/When，而是给每一步产生结构化语义状态。
  \item \textbf{把 RNNRepair 的状态抽象迁移到 LLM agent 日志。} 原来 RNNRepair 聚类 RNN hidden state；我们聚类 LLM semantic profile。
  \item \textbf{每个 agent 建自己的状态机。} 这比全局聚类更符合多智能体系统的角色差异。
  \item \textbf{用状态风险表做可解释归因。} 每个预测都能回看：它是因为 profile 可疑，还是因为历史上同类状态风险高。
  \item \textbf{支持缓存。} LLM profiling 可以缓存；后续改 GMM、risk table、decoding 策略时，不必重复调用 LLM。
\end{enumerate}

\section{目前还需要补强的地方}

如果要往论文级别推进，建议下一步补这些实验：

\begin{enumerate}
  \item \textbf{调参要更严格。} 当前 \(\lambda=2\) 和 earliest 阈值 4 是经验值，应该在 validation set 上调，或者做 nested CV。
  \item \textbf{做 ablation。} 比较 only profile、only risk、combined、不同 window、不同 K 选择策略。
  \item \textbf{补 algorithm-generated 子集。} 当前完整跑通的是 hand-crafted，需要看算法生成子集上是否也有效。
  \item \textbf{补 DeepSeek baseline。} 本地目前只找到 DeepSeek-v4-flash all-at-once 结果，还缺 step-by-step 和 binary-search 的同模型复现。
  \item \textbf{增强文本特征。} 当前 \texttt{agent\_intent} 和 \texttt{state\_summary} 没进模型，后续可以 embedding 后加入状态表示。
\end{enumerate}

\section{一句话版本}

我们的方法是：先让 LLM 逐步阅读日志，把每一步翻译成 ``动作类型、证据质量、是否偏题、是否有 unsupported claim、工具是否错误'' 这类结构化语义标签；然后像 RNNRepair 一样，把同一个 agent 的语义标签聚成少数几个状态；再从训练日志里统计哪些状态和状态转移最容易对应决定性错误；最后在新日志中逐步打分，选择最可疑的 agent 和最早错误 step。

\end{document}
